\documentclass[12pt,a4paper]{article}

\usepackage[spanish,es-nodecimaldot]{babel}
\usepackage[utf8]{inputenc}
\usepackage[T1]{fontenc}
\usepackage{lmodern}
\usepackage{geometry}
\usepackage{microtype}
\usepackage{setspace}
\usepackage{parskip}
\usepackage{enumitem}
\usepackage{booktabs}
\usepackage{xcolor}
\usepackage{listings}
\usepackage{hyperref}

\geometry{margin=2.5cm}
\onehalfspacing
\setlist[itemize]{leftmargin=*,itemsep=0.2em}
\setlist[enumerate]{leftmargin=*,itemsep=0.2em}

\hypersetup{
  colorlinks=true,
  linkcolor=black,
  urlcolor=blue,
  pdftitle={Resolución del Práctico 1 - Testing de Software 2024},
  pdfauthor={}
}

\lstdefinestyle{javastyle}{
  language=Java,
  basicstyle=\ttfamily\small,
  keywordstyle=\color{blue!60!black},
  commentstyle=\color{green!40!black},
  stringstyle=\color{orange!50!black},
  numbers=left,
  numberstyle=\tiny,
  stepnumber=1,
  numbersep=8pt,
  frame=single,
  breaklines=true,
  tabsize=2,
  showstringspaces=false
}

\title{\textbf{Resolución del Práctico 1}\\Testing de Software 2024}
\author{}
\date{}

\begin{document}

\maketitle
\tableofcontents
\newpage

\section*{Introducción}
\addcontentsline{toc}{section}{Introducción}
El presente documento unifica la resolución del \textit{Práctico 1} de la materia \textit{Testing de Software 2024}. Se preserva el contenido técnico original de cada ejercicio y se mejora su presentación mediante una redacción académica, una estructura homogénea y una notación consistente.

\section{Ejercicio 1}
Lectura de los capítulos 1 y 2 del libro \textit{Introduction to Software Testing}.

\section{Ejercicio 2}
\subsection{Diferencia entre \textit{fault} y \textit{failure}}
La diferencia central es la siguiente:
\begin{itemize}
  \item \textbf{Fault (defecto):} problema interno en el software (código, diseño o especificación).
  \item \textbf{Failure (falla):} comportamiento externo incorrecto, observable durante la ejecución.
\end{itemize}

Además, entre ambos conceptos aparece \textbf{error}, entendido como un estado interno incorrecto que surge al ejecutar un defecto.

En términos causales:
\[
\textit{fault} \rightarrow \textit{error} \rightarrow \textit{failure}
\]

Por lo tanto, un defecto puede existir incluso sin generar una falla visible en toda ejecución: para que la falla ocurra, el estado erróneo debe propagarse hasta una salida observable.

\section{Ejercicio 3}
Se analizaron cuatro programas defectuosos. Para cada caso se identifica el defecto, una corrección posible y los casos pedidos en los incisos b), c) y d).

\subsection{Programa 1: \texttt{findLast(int[] x, int y)}}
\subsubsection*{a) Defecto y corrección}
El defecto está en la condición del ciclo:
\begin{lstlisting}[style=javastyle]
for (int i = x.length - 1; i > 0; i--)
\end{lstlisting}
De esta forma nunca se evalúa el índice \texttt{0}. La corrección es:
\begin{lstlisting}[style=javastyle]
for (int i = x.length - 1; i >= 0; i--)
\end{lstlisting}

\subsubsection*{b) Caso que no ejecute el defecto}
No es posible con entradas válidas: la condición defectuosa del \texttt{for} se evalúa en toda invocación del método.

\subsubsection*{c) Caso que ejecute el defecto y no produzca falla}
\texttt{x = [5,2,3], y = 2} \hspace{0.5cm} Resultado esperado: \texttt{1}.

\subsubsection*{d) Caso que ejecute el defecto y produzca falla}
\texttt{x = [2,3,5], y = 2} \hspace{0.5cm} Esperado: \texttt{0}. Resultado defectuoso: \texttt{-1}.

\subsection{Programa 2: \texttt{lastZero(int[] x)}}
\subsubsection*{a) Defecto y corrección}
El método debe devolver el \textit{último} índice con valor cero, pero recorre de izquierda a derecha y retorna el primer cero encontrado. La corrección consiste en recorrer desde el final.
\begin{lstlisting}[style=javastyle]
for (int i = x.length - 1; i >= 0; i--) {
    if (x[i] == 0) {
        return i;
    }
}
\end{lstlisting}

\subsubsection*{b) Caso que no ejecute el defecto}
No es posible para arreglos no nulos, ya que la estrategia defectuosa forma parte del recorrido principal.

\subsubsection*{c) Caso que ejecute el defecto y no produzca falla}
\texttt{x = [1,0,2]} \hspace{0.5cm} Resultado esperado: \texttt{1}.

\subsubsection*{d) Caso que ejecute el defecto y produzca falla}
\texttt{x = [0,1,0]} \hspace{0.5cm} Esperado: \texttt{2}. Resultado defectuoso: \texttt{0}.

\subsection{Programa 3: \texttt{countPositive(int[] x)}}
\subsubsection*{a) Defecto y corrección}
Se cuenta \texttt{x[i] >= 0}, lo que incluye al cero como positivo. La condición correcta es:
\begin{lstlisting}[style=javastyle]
if (x[i] > 0) {
    count++;
}
\end{lstlisting}

\subsubsection*{b) Caso que no ejecute el defecto}
\texttt{x = []}. El cuerpo del \texttt{for} no se ejecuta.

\subsubsection*{c) Caso que ejecute el defecto y no produzca falla}
\texttt{x = [-4,2,2]} \hspace{0.5cm} Resultado esperado: \texttt{2}.

\subsubsection*{d) Caso que ejecute el defecto y produzca falla}
\texttt{x = [-4,2,0,2]} \hspace{0.5cm} Esperado: \texttt{2}. Resultado defectuoso: \texttt{3}.

\subsection{Programa 4: \texttt{oddOrPos(int[] x)}}
\subsubsection*{a) Defecto y corrección}
La implementación defectuosa usa \texttt{x[i] \% 2 == 1}. En Java, los impares negativos no satisfacen esa condición (por ejemplo, \texttt{-3 \% 2 == -1}).

La corrección es:
\begin{lstlisting}[style=javastyle]
if (x[i] % 2 != 0 || x[i] > 0) {
    count++;
}
\end{lstlisting}

\subsubsection*{b) Caso que no ejecute el defecto}
\texttt{x = []}. El cuerpo del \texttt{for} no se ejecuta.

\subsubsection*{c) Caso que ejecute el defecto y no produzca falla}
\texttt{x = [2,4,-2]} \hspace{0.5cm} Resultado esperado: \texttt{2}.

\subsubsection*{d) Caso que ejecute el defecto y produzca falla}
\texttt{x = [-3,-2,0,1,4]} \hspace{0.5cm} Esperado: \texttt{3}. Resultado defectuoso: \texttt{2}.

\subsection{Síntesis del ejercicio}
\begin{itemize}
  \item \texttt{findLast}: no evalúa el índice \texttt{0}.
  \item \texttt{lastZero}: devuelve el primer cero en vez del último.
  \item \texttt{countPositive}: contabiliza el \texttt{0} como positivo.
  \item \texttt{oddOrPos}: no reconoce impares negativos.
\end{itemize}

\section{Ejercicio 4}
Se implementaron las reparaciones del ejercicio anterior y se incorporaron casos de prueba JUnit para verificar el comportamiento corregido.

\subsection{Reparaciones implementadas}
\begin{itemize}
  \item \texttt{findLast}: cambio de condición de recorrido de \texttt{i > 0} a \texttt{i >= 0}.
  \item \texttt{lastZero}: recorrido desde el final del arreglo.
  \item \texttt{countPositive}: reemplazo de \texttt{>= 0} por \texttt{> 0}.
  \item \texttt{oddOrPos}: reemplazo de \texttt{\% 2 == 1} por \texttt{\% 2 != 0}.
\end{itemize}

\subsection{Casos de prueba JUnit}
Se implementaron los mismos casos del enunciado que evidenciaban fallas en las versiones defectuosas:
\begin{itemize}
  \item \texttt{findLast([2,3,5], 2) = 0}
  \item \texttt{lastZero([0,1,0]) = 2}
  \item \texttt{countPositive([-4,2,0,2]) = 2}
  \item \texttt{oddOrPos([-3,-2,0,1,4]) = 3}
\end{itemize}

\subsection{Detección de falla usando el modelo RIPR}
\subsubsection*{1) \texttt{findLast}}
\begin{itemize}
  \item \textbf{Reachability:} el test alcanza la condición defectuosa del \texttt{for}.
  \item \textbf{Infection:} se omite la revisión de \texttt{x[0]}.
  \item \textbf{Propagation:} se retorna un índice incorrecto.
  \item \textbf{Revealability:} la aserción muestra diferencia entre esperado y obtenido.
\end{itemize}

\subsubsection*{2) \texttt{lastZero}}
\begin{itemize}
  \item \textbf{Reachability:} se ejecuta el recorrido de izquierda a derecha.
  \item \textbf{Infection:} se selecciona un candidato incorrecto (primer cero).
  \item \textbf{Propagation:} ese índice se retorna como resultado final.
  \item \textbf{Revealability:} la aserción detecta \texttt{0} en lugar de \texttt{2}.
\end{itemize}

\subsubsection*{3) \texttt{countPositive}}
\begin{itemize}
  \item \textbf{Reachability:} se evalúa \texttt{x[i] >= 0}.
  \item \textbf{Infection:} el contador se incrementa indebidamente para \texttt{0}.
  \item \textbf{Propagation:} el total final queda sobreestimado.
  \item \textbf{Revealability:} la aserción muestra \texttt{3} en lugar de \texttt{2}.
\end{itemize}

\subsubsection*{4) \texttt{oddOrPos}}
\begin{itemize}
  \item \textbf{Reachability:} se ejecuta la condición lógica del método.
  \item \textbf{Infection:} los impares negativos no se contabilizan.
  \item \textbf{Propagation:} el conteo final queda una unidad por debajo.
  \item \textbf{Revealability:} la aserción detecta \texttt{2} en lugar de \texttt{3}.
\end{itemize}

\subsection{Ejecución de pruebas}
Desde la carpeta del ejercicio:
\begin{center}
\texttt{mvn -Dmaven.repo.local=.m2 test}
\end{center}

\section{Ejercicio 5}
Se analizaron las clases \texttt{Point} y \texttt{ColorPoint} del paquete \texttt{practico1\_exercises.color\_points}, en el contexto del problema clásico de \texttt{equals} en jerarquías por herencia.

\subsection{Análisis de \texttt{Point}}
\subsubsection*{a) Problema y modificación propuesta}
\texttt{Point} define igualdad por valor y permite que subclases también sean tratadas como \texttt{Point} (uso de \texttt{instanceof Point}). Esto habilita inconsistencias cuando una subclase agrega estado adicional (por ejemplo, color).

La solución recomendada es evitar extender \texttt{Point} para agregar estado de valor: usar composición (o, alternativamente, declarar \texttt{Point} como \texttt{final}).

\subsubsection*{b) Test que no ejecute la falla}
Comparación entre dos instancias de \texttt{Point} con mismas coordenadas.

\subsubsection*{c) Test que ejecute la falla sin error observable}
Comparación \texttt{Point(1,2)} con \texttt{ColorPoint(2,3,RED)}. Aunque se ejecuta la lógica problemática de comparación entre tipos, el resultado sigue siendo correcto (\texttt{false}) por diferencia de coordenadas.

\subsubsection*{d) Test con error pero sin falla}
No se considera posible en este caso: cuando aparece el error lógico en \texttt{equals}, se manifiesta directamente como comportamiento observable incorrecto.

\subsection{Análisis de \texttt{ColorPoint}}
\subsubsection*{a) Problema y modificación propuesta}
\texttt{ColorPoint} redefine \texttt{equals} exigiendo \texttt{instanceof ColorPoint}, lo que rompe simetría con \texttt{Point}. Por ejemplo, puede ocurrir:
\begin{center}
\texttt{p.equals(cp) == true} y \texttt{cp.equals(p) == false}
\end{center}

La corrección aplicada consiste en usar composición en lugar de herencia (implementada en \texttt{ColorPointFixed}).

\subsubsection*{b) Test que no ejecute la falla}
Comparación entre dos \texttt{ColorPoint} con mismas coordenadas y color.

\subsubsection*{c) Test que ejecute la falla sin error observable}
Comparación \texttt{ColorPoint(1,2,RED)} con \texttt{Point(4,5)}. Se ejecuta la rama conflictiva, pero el resultado esperado sigue siendo \texttt{false}.

\subsubsection*{d) Test con error pero sin falla}
No se considera posible, por la misma razón conceptual indicada para \texttt{Point}.

\refstepcounter{section}
\section*{Ejercicio \thesection}
\addcontentsline{toc}{section}{Ejercicio \thesection}
\subsection{Aclaración sobre el material provisto}
En el material disponible no estaban provistas la clase \texttt{PointTest} ni el paquete \texttt{practico1\_exercises.point\_set}. Por ello se crearon:
\begin{itemize}
  \item una implementación de \texttt{Point},
  \item una clase de pruebas \texttt{PointTest},
  \item tests adicionales para cubrir comportamiento en \texttt{HashSet}.
\end{itemize}

\subsection{a) Nuevos tests, falla detectada y corrección}
Se incorporaron pruebas para:
\begin{itemize}
  \item igualdad entre puntos equivalentes,
  \item desigualdad entre puntos distintos,
  \item tratamiento de duplicados en \texttt{HashSet},
  \item búsqueda de elementos equivalentes en \texttt{HashSet}.
\end{itemize}

La falla detectada fue la inconsistencia entre \texttt{equals} y \texttt{hashCode}. La corrección aplicada fue:
\begin{lstlisting}[style=javastyle]
@Override
public int hashCode() {
    return Objects.hash(x, y);
}
\end{lstlisting}

\subsection{b) Identificación de Arrange--Act--Assert}
En cada test se dejaron explícitas las tres fases:
\begin{itemize}
  \item \textbf{Arrange:} preparación del escenario y datos.
  \item \textbf{Act:} ejecución de la operación bajo prueba.
  \item \textbf{Assert:} verificación de resultados esperados.
\end{itemize}

\subsection{Ejecución de pruebas}
\begin{center}
\texttt{mvn -Dmaven.repo.local=.m2 test}
\end{center}

\section{Ejercicio 7}
\subsection{Aclaración sobre el material provisto}
No se encontraron implementaciones dadas de \texttt{PointSet} ni \texttt{PointSetTest}. Para resolver el ejercicio se crearon ambas clases y se diseñaron pruebas adicionales.

\subsection{Tests implementados}
Se incorporaron casos para verificar:
\begin{itemize}
  \item estado inicial vacío,
  \item no inserción de duplicados equivalentes,
  \item reconocimiento de equivalencia en \texttt{contains},
  \item eliminación de equivalentes en \texttt{remove},
  \item respuesta negativa de \texttt{contains} ante ausencias reales.
\end{itemize}

\subsection{Resultado y análisis}
Los tests pasan sobre la implementación actual. No se detectaron fallas adicionales en \texttt{PointSet}. Este resultado es consistente con la corrección previa de \texttt{Point} (coherencia entre \texttt{equals} y \texttt{hashCode}), dado que \texttt{PointSet} se apoya en un \texttt{HashSet<Point>}.

\subsection{Identificación de Arrange--Act--Assert}
Al igual que en el ejercicio 6, todos los tests se estructuraron con fases explícitas de \textit{arrange}, \textit{act} y \textit{assert}.

\section*{Conclusiones}
\addcontentsline{toc}{section}{Conclusiones}
La resolución del Práctico 1 permitió:
\begin{itemize}
  \item distinguir formalmente los conceptos de \textit{fault}, \textit{error} y \textit{failure};
  \item analizar defectos concretos en programas imperativos y su reparación;
  \item aplicar el modelo RIPR para explicar detección de fallas;
  \item abordar un problema clásico de diseño orientado a objetos en el contrato de \texttt{equals};
  \item reforzar diseño de pruebas unitarias y consistencia semántica en estructuras de datos basadas en hashing.
\end{itemize}

\end{document}
